\documentclass[aspectratio=169]{beamer}

\usepackage[utf8]{inputenc}
\usepackage[T1]{fontenc}
\usepackage[french]{babel}
\usepackage{graphicx}
\usepackage{booktabs}
\usepackage{xcolor}
\usepackage{amsmath}
\usepackage{array}
\usepackage{colortbl}
\usepackage{listings}

% ── Thème ──────────────────────────────────────────────────────
\usetheme{Madrid}
\definecolor{medblue}{RGB}{10, 75, 140}
\definecolor{medlight}{RGB}{220, 235, 255}
\definecolor{codegray}{RGB}{245, 245, 245}
\setbeamercolor{structure}{fg=medblue}
\setbeamercolor{block title}{bg=medblue, fg=white}
\setbeamercolor{block body}{bg=medlight, fg=black}
\setbeamercolor{alerted text}{fg=red!70!black}

\setbeamersize{text margin left=6mm,text margin right=6mm}
\setbeamertemplate{navigation symbols}{}
\setlength{\abovedisplayskip}{2pt}
\setlength{\belowdisplayskip}{2pt}
\setlength{\abovedisplayshortskip}{1pt}
\setlength{\belowdisplayshortskip}{1pt}

\lstset{
  basicstyle=\ttfamily\tiny,
  backgroundcolor=\color{codegray},
  frame=single, framerule=0pt,
  breaklines=true,
  keywordstyle=\color{medblue}\bfseries,
  commentstyle=\color{gray},
  language=Python,
  showstringspaces=false,
}

\newcommand{\good}[1]{\textcolor{green!55!black}{\textbf{#1}}}
\newcommand{\warn}[1]{\textcolor{orange!80!black}{\textbf{#1}}}
\newcommand{\bad}[1]{\textcolor{red!70!black}{\textbf{#1}}}

% ── Métadonnées ────────────────────────────────────────────────
\title[CNN -- Détection du Mélanome]
      {Classification d'Images Médicales\\par Réseaux de Neurones Convolutifs}
\subtitle{Application à la Détection du Mélanome Cutané}
\author[Aitelamiri, NOM]{Amine AITELAMIRI \and [Prénom NOM]}
\institute{[Établissement \textbullet{} Filière]}
\date{Mai 2026}

\begin{document}

% ============================================================
% Slide titre
% ============================================================
\begin{frame}
  \titlepage
\end{frame}

% ============================================================
% 1. Le dataset
% ============================================================
\begin{frame}{1. Le dataset}
  \begin{columns}[T,onlytextwidth]
    \column{0.46\textwidth}
      \begin{block}{Problème médical}
        \small Classification \textbf{binaire} : Bénin / Malin\\[2pt]
        Détection précoce : survie à 5 ans de \bad{25\%} à \good{99\%}
      \end{block}
      \vspace{3pt}
      \begin{block}{Distribution}
        \small\renewcommand{\arraystretch}{1.1}
        \begin{tabular}{lcc}
          \toprule & \textbf{Bénin} & \textbf{Malin} \\ \midrule
          Train      & 6\,289 & 5\,590 \\
          Validation & 1\,000 & 1\,000 \\ \bottomrule
        \end{tabular}\\[2pt]
        $224\times224$ px, RGB --- Quasi-équilibré
      \end{block}
      \vspace{3pt}
      \begin{alertblock}{Difficultés anticipées}
        \small Lésions visuellement similaires,
        risque d'overfitting (dataset de taille modérée)
      \end{alertblock}
    \column{0.51\textwidth}
      \centering
      \includegraphics[width=\linewidth,height=0.72\textheight,keepaspectratio]{dataset_examples.png}
  \end{columns}
\end{frame}

% ============================================================
% 2. Représentation des données
% ============================================================
\begin{frame}[fragile]{2. Représentation des données --- tenseurs et canaux RGB}
  \begin{columns}[T,onlytextwidth]
    \column{0.53\textwidth}
      \centering
      \includegraphics[width=\linewidth,height=0.68\textheight,keepaspectratio]{image_tenseur.png}
    \column{0.44\textwidth}
      \begin{block}{Une image = tenseur \texttt{[C,H,W]}}
        \small
        \begin{itemize}\setlength\itemsep{4pt}
          \item Forme : $[3, 224, 224]$ (RGB)
          \item \texttt{ToTensor()} : pixels [0--255] $\to$ float [0,1]
          \item Pixel PIL : \texttt{(181, 140, 144)}\\
                Pixel tenseur : \texttt{0.710, 0.549, 0.565}
        \end{itemize}
      \end{block}
      \vspace{4pt}
      \begin{block}{Code}
\begin{lstlisting}
img = Image.open(path).convert("RGB")
t = transforms.ToTensor()(img)
# t.shape: [3, 224, 224]
# min: 0.0000  max: 0.9412
\end{lstlisting}
      \end{block}
  \end{columns}
\end{frame}

% ============================================================
% 3. Dataset et DataLoader
% ============================================================
\begin{frame}[fragile]{3. Dataset et DataLoader}
  \begin{columns}[T,onlytextwidth]
    \column{0.49\textwidth}
      \centering
      \includegraphics[width=\linewidth,height=0.50\textheight,keepaspectratio]{batch_train.png}\\
      \tiny Batch \texttt{[32, 3, 224, 224]}
    \column{0.48\textwidth}
      \begin{block}{\texttt{MelonaDataset}}
\begin{lstlisting}
class MelonaDataset(Dataset):
    def __init__(self, root, transform):
        # sous-dossiers -> classes
        # liste de (chemin, label)
    def __getitem__(self, idx):
        img = Image.open(path)
        return transform(img), label
\end{lstlisting}
      \end{block}
      \vspace{2pt}
      \begin{block}{\texttt{DataLoader}}
\begin{lstlisting}
loader = DataLoader(
    dataset, batch_size=32,
    shuffle=True, num_workers=0)
\end{lstlisting}
        \small \textbf{Shuffle} : mélange chaque époque.\\
        \textbf{Batch} : 32 images en parallèle sur GPU.
      \end{block}
  \end{columns}
\end{frame}

% ============================================================
% 4. Comprendre l'apprentissage
% ============================================================
\begin{frame}[fragile]{4. Comprendre l'apprentissage}
  \begin{columns}[T,onlytextwidth]
    \column{0.48\textwidth}
      \begin{block}{Concepts clés}
        \small
        \begin{description}\setlength\itemsep{5pt}
          \item[\textbf{Loss}] CrossEntropyLoss :
                $\mathcal{L} = -\sum_c y_c \log \hat{p}_c$
                (softmax implicite)
          \item[\textbf{Gradient}] $\nabla_\theta \mathcal{L}$
                par rétropropagation (\texttt{backward()})
          \item[\textbf{Époque}] Passage complet sur le dataset.
                Entraînement = 20 époques.
          \item[\textbf{Accuracy}] corrects\,/\,total
        \end{description}
      \end{block}
    \column{0.48\textwidth}
      \begin{block}{Boucle d'entraînement}
\begin{lstlisting}
for imgs, lbls in train_loader:
    optimizer.zero_grad()
    out  = model(imgs.to(device))
    loss = criterion(out, lbls.to(device))
    loss.backward()
    optimizer.step()
\end{lstlisting}
      \end{block}
      \vspace{3pt}
      \begin{block}{Évaluation}
\begin{lstlisting}
model.eval()            # desactive Dropout
with torch.no_grad():   # pas de gradient
    preds = model(imgs).argmax(dim=1)
\end{lstlisting}
      \end{block}
  \end{columns}
\end{frame}

% ============================================================
% 5. Premier CNN
% ============================================================
\begin{frame}{5. Premier CNN --- architecture et courbes}
  \begin{columns}[T,onlytextwidth]
    \column{0.37\textwidth}
      \begin{block}{SimpleCNN}
        \tiny\renewcommand{\arraystretch}{1.1}
        \begin{tabular}{>{\bfseries}lc}
          \toprule Couche & Sortie \\ \midrule
          Input              & $3\times224\times224$ \\
          Conv(3$\to$16)+Pool  & $16\times112\times112$ \\
          Conv(16$\to$32)+Pool & $32\times56\times56$  \\
          Conv(32$\to$64)+Pool & $64\times28\times28$  \\
          Flatten            & $50\,176$             \\
          Linear+ReLU        & $256$                 \\
          Dropout(0.5)       & $256$                 \\
          Linear             & $2$                   \\ \bottomrule
        \end{tabular}\\[3pt]
        \small \textbf{12.9M} paramètres\\
        Adam lr=$10^{-3}$, 20 époques
      \end{block}
    \column{0.59\textwidth}
      \centering
      \includegraphics[width=\linewidth,height=0.44\textheight,keepaspectratio]{courbes_cnn_simple.png}
      \vspace{2pt}
      \begin{alertblock}{Limites observées}
        \small Gap \warn{94.7\%} train vs \warn{91.4\%} val ($\Delta=3.3\%$) --- overfitting naissant.
        Loss val remonte : le modèle mémorise.
      \end{alertblock}
  \end{columns}
\end{frame}

% ============================================================
% 6. Normalisation
% ============================================================
\begin{frame}[fragile]{6. Normalisation --- pourquoi, comment, impact}
  \begin{columns}[T,onlytextwidth]
    \column{0.53\textwidth}
      \centering
      \includegraphics[width=\linewidth,height=0.70\textheight,keepaspectratio]{courbes_cnn_normalise.png}
    \column{0.44\textwidth}
      \begin{block}{Pourquoi ?}
        \small Distribution centrée sur 0.72 (peau claire) :
        gradients déséquilibrés. Normaliser centre et réduit chaque canal.
      \end{block}
      \vspace{3pt}
      \begin{block}{Comment (stats sur TRAIN)}
\begin{lstlisting}
mean = [0.723, 0.556, 0.539]
std  = [0.183, 0.193, 0.208]
Normalize(mean=mean, std=std)
\end{lstlisting}
      \end{block}
      \vspace{3pt}
      \begin{block}{Impact}
        \small Val : 91.4\% $\to$ \good{93.3\%} (\textbf{+1.9\%})\\
        Gap : 3.3\% $\to$ \bad{5.3\%}\\
        \textit{Overfitting aggravé : convergence meilleure,\\
        mais le modèle généralise moins.}
      \end{block}
  \end{columns}
\end{frame}

% ============================================================
% 7. Data augmentation
% ============================================================
\begin{frame}{7. Data augmentation --- stratégie et impact}
  \begin{columns}[T,onlytextwidth]
    \column{0.53\textwidth}
      \centering
      \includegraphics[width=\linewidth,height=0.70\textheight,keepaspectratio]{courbes_cnn_aug.png}
    \column{0.44\textwidth}
      \begin{block}{Stratégie (train uniquement)}
        \small
        \begin{itemize}\setlength\itemsep{4pt}
          \item \texttt{RandomHorizontalFlip}$(p{=}0.5)$
          \item \texttt{RandomRotation}$(\pm10°)$
          \item \texttt{ColorJitter}(brightness, contrast = 0.2)
        \end{itemize}
      \end{block}
      \vspace{3pt}
      \begin{block}{Impact}
        \small\renewcommand{\arraystretch}{1.1}
        \begin{tabular}{lcc}
          \toprule Config & Val & Gap \\ \midrule
          Norm. seule  & 93.3\% & \bad{5.3\%} \\
          + Augment.   & \good{91.5\%} & \good{$\approx 0$\%} \\
          \bottomrule
        \end{tabular}\\[3pt]
        Gap fermé : le modèle \textbf{généralise}.
      \end{block}
  \end{columns}
\end{frame}

% ============================================================
% 8. Optimiseur et learning rate
% ============================================================
\begin{frame}{8. Choix de l'optimiseur et du learning rate}
  \begin{columns}[T,onlytextwidth]
    \column{0.53\textwidth}
      \centering
      \includegraphics[width=\linewidth,height=0.70\textheight,keepaspectratio]{courbes_optimiseurs.png}
    \column{0.44\textwidth}
      \begin{block}{Résultats}
        \small\renewcommand{\arraystretch}{1.2}
        \begin{tabular}{lc}
          \toprule Optimiseur & Acc. val \\ \midrule
          Adam $10^{-3}$ & 91.3\% \\
          Adam $10^{-4}$ & \good{91.5\%} \\
          SGD  $10^{-2}$ & \warn{89.0\%} \\ \bottomrule
        \end{tabular}
      \end{block}
      \vspace{3pt}
      \begin{block}{Pourquoi Adam ?}
        \small LR adaptatif par paramètre --- robuste sur images hétérogènes.
        SGD plus bruité, sensible au réglage du LR.
      \end{block}
      \vspace{3pt}
      \begin{alertblock}{Retenu}
        \small Adam $10^{-3}$ (compromis vitesse / qualité)
      \end{alertblock}
  \end{columns}
\end{frame}

% ============================================================
% 9. Scheduler
% ============================================================
\begin{frame}{9. Learning Rate Scheduler --- StepLR}
  \begin{columns}[T,onlytextwidth]
    \column{0.52\textwidth}
      \centering
      \includegraphics[width=0.95\linewidth,height=0.36\textheight,keepaspectratio]{courbes_scheduler.png}\\[3pt]
      \includegraphics[width=0.95\linewidth,height=0.36\textheight,keepaspectratio]{courbes_lr_schedule.png}
    \column{0.45\textwidth}
      \begin{block}{Principe --- step=7, $\gamma=0.1$}
        \small
        \begin{tabular}{lc}
          \toprule Époques & LR \\ \midrule
          1--7   & $10^{-3}$ \\
          8--14  & $10^{-4}$ \\
          15--20 & $10^{-5}$ \\ \bottomrule
        \end{tabular}
      \end{block}
      \vspace{4pt}
      \begin{block}{Impact}
        \small Sans scheduler : 91.0\% val\\
        Avec StepLR : \good{92.0\%}\\[2pt]
        Convergence plus stable après chaque palier.
        Le graphe log confirme les 2 sauts.
      \end{block}
  \end{columns}
\end{frame}

% ============================================================
% 10. TensorBoard
% ============================================================
\begin{frame}[fragile]{10. Suivi des expériences --- TensorBoard}
  \begin{columns}[T,onlytextwidth]
    \column{0.53\textwidth}
      \centering
      \includegraphics[width=\linewidth,height=0.62\textheight,keepaspectratio]{courbes_optimiseurs.png}\\
      \small\textit{Approche équivalente : matplotlib superposé}
    \column{0.44\textwidth}
      \begin{block}{TensorBoard}
\begin{lstlisting}
from torch.utils.tensorboard \
    import SummaryWriter
w = SummaryWriter("runs/adam_1e3")
w.add_scalar("Loss/val", loss, ep)
w.add_scalar("Acc/val",  acc,  ep)
# tensorboard --logdir runs/
\end{lstlisting}
        \small Comparaison interactive dans le navigateur,
        zoom, histogrammes de poids.
      \end{block}
  \end{columns}
\end{frame}

% ============================================================
% 11. ResNet — transfer learning, backbone gelé
% ============================================================
\begin{frame}[fragile]{11. ResNet18 --- Transfer Learning}
  \begin{columns}[T,onlytextwidth]
    \column{0.48\textwidth}
      \begin{block}{Backbone gelé}
\begin{lstlisting}
resnet = models.resnet18(
    weights=IMAGENET1K_V1)
for p in resnet.parameters():
    p.requires_grad = False  # geler tout
resnet.fc = nn.Linear(512, num_classes)
# 1 026 params entrainables seulement
\end{lstlisting}
        \small Norm. ImageNet : mean=[0.485, 0.456, 0.406]
      \end{block}
      \vspace{3pt}
      \begin{alertblock}{Résultat backbone gelé}
        \small Acc. val = \warn{88.9\%} --- \textit{inférieur au SimpleCNN !}\\
        Features ImageNet $\neq$ patterns dermoscopiques.
      \end{alertblock}
    \column{0.48\textwidth}
      \centering
      \includegraphics[width=\linewidth,height=0.68\textheight,keepaspectratio]{courbes_resnet.png}\\
      \small\textit{Courbe bleue : backbone gelé}
  \end{columns}
\end{frame}

% ============================================================
% 12. Fine-tuning
% ============================================================
\begin{frame}[fragile]{12. Fine-tuning --- dégelage partiel de layer4}
  \begin{columns}[T,onlytextwidth]
    \column{0.50\textwidth}
      \begin{block}{Stratégie}
\begin{lstlisting}
for p in resnet.layer4.parameters():
    p.requires_grad = True   # degeler

optimizer = Adam([
  {"params": resnet.layer4.parameters(),
   "lr": 1e-4},  # adaptation douce
  {"params": resnet.fc.parameters(),
   "lr": 1e-3},  # apprentissage rapide
])
# 8.4M entrainables / 11.2M total
\end{lstlisting}
      \end{block}
      \vspace{3pt}
      \begin{block}{Gain}
        \small\renewcommand{\arraystretch}{1.15}
        \begin{tabular}{lc}
          \toprule Modèle & Acc. val \\ \midrule
          Backbone gelé      & \warn{88.9\%} \\
          SimpleCNN + norm.  & 93.3\% \\
          Fine-tuning layer4 & \good{\textbf{95.6\%}} \\ \bottomrule
        \end{tabular}
      \end{block}
    \column{0.46\textwidth}
      \centering
      \includegraphics[width=\linewidth,height=0.38\textheight,keepaspectratio]{courbes_resnet.png}
      \vspace{4pt}
      \begin{alertblock}{Pourquoi layer4 ?}
        \small Layer4 = features de haut niveau spécifiques au mélanome.
        Couches basses (bords, contours) = universelles.
        Évite le \textit{catastrophic forgetting}.
      \end{alertblock}
  \end{columns}
\end{frame}

% ============================================================
% 13. Analyse des erreurs
% ============================================================
\begin{frame}{13. Analyse des erreurs --- Matrice de confusion}
  \begin{columns}[T,onlytextwidth]
    \column{0.40\textwidth}
      \centering
      \includegraphics[width=\linewidth,height=0.50\textheight,keepaspectratio]{matrice_confusion.png}\\[2pt]
      \includegraphics[width=\linewidth,height=0.22\textheight,keepaspectratio]{exemples_incorrects.png}\\
      \tiny\textit{Exemples mal classés}
    \column{0.56\textwidth}
      \begin{block}{Rapport de classification (2\,000 images)}
        \small\renewcommand{\arraystretch}{1.1}
        \begin{tabular}{lcccc}
          \toprule Classe & Préc. & Rappel & F1 & N \\ \midrule
          Bénin & 95\% & 96\% & 96\% & 1\,000 \\
          Malin & 96\% & 95\% & 96\% & 1\,000 \\ \midrule
          \multicolumn{2}{l}{\textbf{Accuracy}} &
            \multicolumn{3}{c}{\good{\textbf{96\% (1\,911 / 2\,000)}}} \\
          \bottomrule
        \end{tabular}
      \end{block}
      \vspace{3pt}
      \begin{alertblock}{Enjeux médicaux}
        \small
        \begin{itemize}\setlength\itemsep{3pt}
          \item \bad{53 faux négatifs} : malins classés bénins --- \textbf{risque vital}
          \item \warn{36 faux positifs} : bénins classés malins --- sur-diagnostic
        \end{itemize}
        \vspace{2pt}
        En clinique : abaisser le seuil de décision
        pour maximiser le \textbf{rappel Malin}.
      \end{alertblock}
  \end{columns}
\end{frame}

% ============================================================
% 14. Conclusion
% ============================================================
\begin{frame}{14. Conclusion --- Quel modèle retenir ?}
  \begin{columns}[T,onlytextwidth]
    \column{0.54\textwidth}
      \begin{block}{Bilan des expériences}
        \tiny\renewcommand{\arraystretch}{1.15}
        \begin{tabular}{clcc}
          \toprule \textbf{P.} & \textbf{Configuration} & \textbf{Train} & \textbf{Val} \\ \midrule
          5  & SimpleCNN (sans norm.)  & 94.7\% & 91.4\% \\
          6  & + Normalisation         & 98.6\% & 93.3\% \\
          7  & + Augmentation          & 90.3\% & 91.5\% \\
          8  & + Adam $10^{-4}$        & 89.4\% & 91.5\% \\
          9  & + StepLR(7, 0.1)        & 90.0\% & 92.0\% \\
          11 & ResNet backbone gelé    & ---    & \warn{88.9\%} \\
          12 & ResNet fine-tuning      & $>$99\%& \good{\textbf{95.6\%}} \\ \midrule
          13 & \textbf{Analyse FT} &
            \multicolumn{2}{c}{\good{\textbf{96\% --- 89 err. / 2\,000}}} \\
          \bottomrule
        \end{tabular}
      \end{block}
    \column{0.42\textwidth}
      \begin{block}{Ce qui a le plus amélioré}
        \small
        \begin{enumerate}\setlength\itemsep{4pt}
          \item \textbf{Fine-tuning} \good{(+4.1\%)}
          \item \textbf{Normalisation} (+1.9\%)
          \item \textbf{Scheduler} (stabilité)
          \item \textbf{Augmentation} (gap $\to$ 0)
        \end{enumerate}
      \end{block}
      \vspace{3pt}
      \begin{block}{Avec plus de temps\ldots}
        \small Early stopping, EfficientNet, Focal loss (rappel Malin), Grad-CAM
      \end{block}
      \vspace{4pt}
      \begin{alertblock}{}
        \centering \large\textbf{Merci} \normalsize --- \textit{Questions ?}
      \end{alertblock}
  \end{columns}
\end{frame}

\end{document}
